%\makeglossaries
\makenoidxglossaries

% Danh mục thuật ngữ và từ viết tắt
\newglossaryentry{cnn}{
    type=\acronymtype,
    name={CNN},
    description={Mạng nơ-ron tích chập (Convolutional Neural Network)},
    first={Mạng nơ-ron tích chập (Convolutional Neural Network - CNN)}
}
\newglossaryentry{lstm}{
    type=\acronymtype,
    name={LSTM},
    description={Mạng bộ nhớ dài-ngắn hạn (Long Short-Term Memory)},
    first={Mạng bộ nhớ dài-ngắn hạn (Long Short-Term Memory - LSTM)}
}
\newglossaryentry{imu}{
    type=\acronymtype,
    name={IMU},
    description={Khối đo lường quán tính (Inertial Measurement Unit)},
    first={Khối đo lường quán tính (Inertial Measurement Unit - IMU)}
}
\newglossaryentry{far}{
    type=\acronymtype,
    name={FAR},
    description={Tỉ lệ chấp nhận sai (False Acceptance Rate)},
    first={Tỉ lệ chấp nhận sai (False Acceptance Rate - FAR)}
}
\newglossaryentry{frr}{
    type=\acronymtype,
    name={FRR},
    description={Tỉ lệ từ chối sai (False Rejection Rate)},
    first={Tỉ lệ từ chối sai (False Rejection Rate - FRR)}
}
\newglossaryentry{eer}{
    type=\acronymtype,
    name={EER},
    description={Tỉ lệ lỗi cân bằng (Equal Error Rate)},
    first={Tỉ lệ lỗi cân bằng (Equal Error Rate - EER)}
}
\newglossaryentry{auc}{
    type=\acronymtype,
    name={AUC},
    description={Diện tích dưới đường cong ROC (Area Under the Curve)},
    first={Diện tích dưới đường cong ROC (Area Under the Curve - AUC)}
}
\newglossaryentry{ewma}{
    type=\acronymtype,
    name={EWMA},
    description={Trung bình động trọng số mũ (Exponentially Weighted Moving Average)},
    first={Trung bình động trọng số mũ (Exponentially Weighted Moving Average - EWMA)}
}
\newglossaryentry{bilstm}{
    type=\acronymtype,
    name={BiLSTM},
    description={Mạng bộ nhớ dài-ngắn hạn hai chiều (Bidirectional Long Short-Term Memory)},
    first={Mạng bộ nhớ dài-ngắn hạn hai chiều (Bidirectional Long Short-Term Memory - BiLSTM)}
}
\newglossaryentry{convlstm}{
    type=\acronymtype,
    name={ConvLSTM},
    description={Mạng LSTM tích chập (Convolutional LSTM)},
    first={Mạng LSTM tích chập (Convolutional LSTM - ConvLSTM)}
}
\newglossaryentry{roc}{
    type=\acronymtype,
    name={ROC},
    description={Đường cong đặc trưng hoạt động của bộ thu nhận (Receiver Operating Characteristic)},
    first={Đường cong đặc trưng hoạt động của bộ thu nhận (Receiver Operating Characteristic - ROC)}
}
\newglossaryentry{relu}{
    type=\acronymtype,
    name={ReLU},
    description={Đơn vị tuyến tính chỉnh lưu (Rectified Linear Unit)},
    first={Đơn vị tuyến tính chỉnh lưu (Rectified Linear Unit - ReLU)}
}
\newglossaryentry{svm}{
    type=\acronymtype,
    name={SVM},
    description={Máy vectơ hỗ trợ (Support Vector Machine)},
    first={Máy vectơ hỗ trợ (Support Vector Machine - SVM)}
}
\newglossaryentry{knn}{
    type=\acronymtype,
    name={KNN},
    description={K láng giềng gần nhất (K-Nearest Neighbors)},
    first={K láng giềng gần nhất (K-Nearest Neighbors - KNN)}
}
\newglossaryentry{dt}{
    type=\acronymtype,
    name={DT},
    description={Cây quyết định (Decision Tree)},
    first={Cây quyết định (Decision Tree - DT)}
}
\newglossaryentry{bn}{
    type=\acronymtype,
    name={BN},
    description={Mạng Bayes (Bayesian Network)},
    first={Mạng Bayes (Bayesian Network - BN)}
}
\newglossaryentry{mlp}{
    type=\acronymtype,
    name={MLP},
    description={Mạng perceptron đa lớp (Multilayer Perceptron)},
    first={Mạng perceptron đa lớp (Multilayer Perceptron - MLP)}
}
\newglossaryentry{qcd}{
    type=\acronymtype,
    name={QCD},
    description={Phát hiện thay đổi nhanh nhất (Quickest Change Detection)},
    first={Phát hiện thay đổi nhanh nhất (Quickest Change Detection - QCD)}
}